% ---- begin paper/tables/a8_vida.tex
\begin{table}[htbp]\centering
\caption{Daily Life: Items and Scales}
\label{tab:a_daily}
\footnotesize\renewcommand{\arraystretch}{1.0}
\begin{threeparttable}
\begin{tabular}{@{}p{68pt}p{212pt}p{150pt}@{}}
\toprule
 & Question, verbatim & Response scale \\
 & (1) & (2) \\
\midrule
\texttt{a8\_a} & Pensando en una semana de lunes a viernes, en promedio ?`a qué hora te despiertas? & Open numeric answer (clock time) \\
\addlinespace[1pt]
\texttt{a8\_b} & Pensando en una semana de lunes a viernes, en promedio ?`a qué hora te duermes? & Open numeric answer (clock time) \\
\addlinespace[1pt]
\texttt{a8\_e} & Pensando en la última semana, ?`acostumbras a quedarte despierto después de acostarte mirando tu celular o Tablet? & 1 = Nunca; 2 = Una vez a la semana; 3 = Más de una vez a la semana; 4 = Todos los días \\
\addlinespace[1pt]
\texttt{a4} & Pensando en un día de lunes a viernes, en promedio ?`Cuánto tiempo al día usas aplicaciones como Instagram, TikTok, WhatsApp, Youtube, Facebook u otras? & 1 = Nada; 2 = Menos de 1 hora; 3 = Entre 1 y menos de 2 horas; 4 = Entre 2 y menos de 3 horas; 5 = Entre 3 y menos de 4 horas; 6 = Entre 4 y menos de 5 horas; 7 = Más de 5 horas. \\
\addlinespace[1pt]
\texttt{a2} & Durante la última semana, ?`cuántos días hiciste deporte o actividad física por al menos 60 minutos? & Open numeric answer (days, 0--7) \\
\addlinespace[1pt]
\texttt{a7} & Pensando en un día de lunes a viernes, en promedio ?`Cuánto tiempo lees libros, revistas o cómics en algún dispositivo electrónico, por ejemplo: Tablet, celular o computador? & 1 = Nada; 2 = Menos de 1 hora; 3 = Entre 1 y menos de 2 horas; 4 = Entre 2 y menos de 3 horas; 5 = Entre 3 y menos de 4 horas; 6 = Entre 4 y menos de 5 horas; 7 = Más de 5 horas. \\
\addlinespace[1pt]
\texttt{a10\_h} & Comida chatarra (completos, hamburguesas, pizzas, papas fritas, etc.)? & Open numeric answer (days, 0--7) \\
\addlinespace[1pt]
\texttt{a11} & Pensando en una semana de lunes a domingo, ?`cuántas veces al día comes frutas? & Open numeric answer (portions per day) \\
\bottomrule
\end{tabular}
\begin{wpnotes}
\textit{Notes:} Question wording and response labels are reproduced verbatim in the original Spanish of the questionnaire, from the variable and value labels of the public microdata; wording the microdata truncates at 80 characters is marked with an ellipsis. Hours of sleep on weekdays are the survey's own computation from the waking and sleeping times (variable tot\_hrs\_sueno\_sem), kept between 3 and 14. Staying up on the phone, time on social media and time reading are standardized; the days of physical activity and of junk food and the daily portions of fruit enter in their own units.
\end{wpnotes}
\end{threeparttable}
\end{table}
% ---- end paper/tables/a8_vida.tex
